
\documentclass{article}
\usepackage{colortbl}
\usepackage{makecell}
\usepackage{multirow}
\usepackage{supertabular}

\begin{document}

\newcounter{utterance}

\centering \large Interaction Transcript for game `chronicle', experiment `discoveries\_hard', episode 2 with Dialogue Architects\_\_qwen3.5{-}9b\_\_SCoRe\_Qwen3.5{-}9B.
\vspace{24pt}

{ \footnotesize  \setcounter{utterance}{1}
\setlength{\tabcolsep}{0pt}
\begin{supertabular}{c@{$\;$}|p{.15\linewidth}@{}p{.15\linewidth}p{.15\linewidth}p{.15\linewidth}p{.15\linewidth}p{.15\linewidth}}
    \# & $\;$A & \multicolumn{4}{c}{Game Master} & $\;\:$B\\
    \hline

    \theutterance \stepcounter{utterance}  
    & & \multicolumn{4}{p{0.6\linewidth}}{
        \cellcolor[rgb]{0.9,0.9,0.9}{
            \makecell[{{p{\linewidth}}}]{
                \texttt{\tiny{[P1$\langle$GM]}}
                \texttt{You are playing a historical deduction game as the NARRATOR.} \\
\\ 
\texttt{You have been given the name of a historical event, discovery, or milestone. Your goal is to help the DETECTIVE figure out which event you are describing — but you must describe it indirectly, without ever naming it.} \\
\\ 
\texttt{Rules:} \\
\texttt{1. Each round, write exactly ONE sentence describing the event. Add NEW information each round — do not repeat what you wrote before.} \\
\texttt{2. You must NEVER use any of the following forbidden keywords (these would make the event too easy to identify): commission, rondon, the rondon commission} \\
\texttt{3. You must NEVER state the event name directly or any close variant of it.} \\
\texttt{4. You may freely describe causes, consequences, the broader context, the sequence of events, and the significance — but only using general, non{-}identifying language.} \\
\\ 
\texttt{Format your response EXACTLY as follows (required):} \\
\texttt{PARAGRAPH: <exactly one sentence describing the event without forbidden keywords>} \\
\\ 
\texttt{The event you must describe (without naming it) is: Rondon Commission} \\
            }
        }
    }
    & \\ \\

    \theutterance \stepcounter{utterance}  
    & \multicolumn{4}{p{0.6\linewidth}}{
        \cellcolor[rgb]{0.9,0.9,0.9}{
            \makecell[{{p{\linewidth}}}]{
                \texttt{\tiny{[P1$\rangle$GM]}}
                \texttt{PARAGRAPH: In the early 20th century, a distinguished Brazilian naturalist and explorer was tasked with leading a massive expedition to survey and demarcate the remote borders of the Amazonian region, a mission that would ultimately transform the area into a protected sanctuary for indigenous peoples and wildlife.} \\
            }
        }
    }
    & & \\ \\

    \theutterance \stepcounter{utterance}  
    & & \multicolumn{4}{p{0.6\linewidth}}{
        \cellcolor[rgb]{0.9,0.9,0.9}{
            \makecell[{{p{\linewidth}}}]{
                \texttt{\tiny{[GM$\rangle$P2]}}
                \texttt{You are playing a historical deduction game as the DETECTIVE.} \\
\\ 
\texttt{A Narrator will describe a historical event, discovery, or milestone to you over several rounds. Each round you receive a new sentence of description that builds on what came before. Your job is to figure out which specific event is being described.} \\
\\ 
\texttt{After EACH new sentence, respond briefly. Keep your entire response under 100 words. Use this exact format (both sections are required):} \\
\\ 
\texttt{ANALYSIS: <2–3 short sentences: what does the growing description suggest? What type of event is this? What are the key features you are noticing?>} \\
\\ 
\texttt{GUESS: <your best guess for the name of the event being described — be specific, e.g. "The French Revolution" not just "a revolution">} \\
\\ 
\texttt{Rules:} \\
\texttt{{-} You may refine your GUESS each round as more information arrives.} \\
\texttt{{-} Your GUESS should name the specific event, not just a broad category.} \\
\texttt{{-} You have up to 10 rounds total.} \\
\texttt{{-} Be concise — quality of reasoning matters more than length.} \\
\\ 
\texttt{The first sentence from the Narrator is coming now.} \\
\\ 
\\ 
\texttt{{[}Sentence 1{]}} \\
\texttt{In the early 20th century, a distinguished Brazilian naturalist and explorer was tasked with leading a massive expedition to survey and demarcate the remote borders of the Amazonian region, a mission that would ultimately transform the area into a protected sanctuary for indigenous peoples and wildlife.} \\
            }
        }
    }
    & \\ \\

    \theutterance \stepcounter{utterance}  
    & & & \multicolumn{4}{p{0.6\linewidth}}{
        \cellcolor[rgb]{0.9,0.9,0.9}{
            \makecell[{{p{\linewidth}}}]{
                \texttt{\tiny{[GM$\langle$P2]}}
                \texttt{ANALYSIS: The description points to a Brazilian expedition in the early 1900s focused on mapping the Amazon and establishing a protected area. The key figures are likely a naturalist and explorer, with the outcome being a sanctuary for indigenous people and wildlife. This suggests a specific conservation milestone rather than a general exploration.} \\
\\ 
\texttt{GUESS: The creation of the Amazon National Forest (Floresta Nacional do Amazonas) or the early establishment of the Amazon Conservation Area.} \\
            }
        }
    }
     \\ \\

    \theutterance \stepcounter{utterance}  
    & & \multicolumn{4}{p{0.6\linewidth}}{
        \cellcolor[rgb]{0.9,0.9,0.9}{
            \makecell[{{p{\linewidth}}}]{
                \texttt{\tiny{[P1$\langle$GM]}}
                \texttt{Round 1 complete. The Detective's current guess: "The creation of the Amazon National Forest (Floresta Nacional do Amazonas) or the early establishment of the Amazon Conservation Area.". Please write a new paragraph with fresh information.} \\
            }
        }
    }
    & \\ \\

    \theutterance \stepcounter{utterance}  
    & \multicolumn{4}{p{0.6\linewidth}}{
        \cellcolor[rgb]{0.9,0.9,0.9}{
            \makecell[{{p{\linewidth}}}]{
                \texttt{\tiny{[P1$\rangle$GM]}}
                \texttt{PARAGRAPH: This monumental undertaking required the construction of a continuous telegraph line stretching over 2,000 kilometers through dense, uncharted wilderness, a feat that not only connected isolated indigenous communities to the outside world but also laid the foundational infrastructure for the region's future conservation efforts.} \\
            }
        }
    }
    & & \\ \\

    \theutterance \stepcounter{utterance}  
    & & \multicolumn{4}{p{0.6\linewidth}}{
        \cellcolor[rgb]{0.9,0.9,0.9}{
            \makecell[{{p{\linewidth}}}]{
                \texttt{\tiny{[GM$\rangle$P2]}}
                \texttt{{[}Sentence 2{]}} \\
\texttt{This monumental undertaking required the construction of a continuous telegraph line stretching over 2,000 kilometers through dense, uncharted wilderness, a feat that not only connected isolated indigenous communities to the outside world but also laid the foundational infrastructure for the region's future conservation efforts.} \\
            }
        }
    }
    & \\ \\

    \theutterance \stepcounter{utterance}  
    & & & \multicolumn{4}{p{0.6\linewidth}}{
        \cellcolor[rgb]{0.9,0.9,0.9}{
            \makecell[{{p{\linewidth}}}]{
                \texttt{\tiny{[GM$\langle$P2]}}
                \texttt{ANALYSIS: The mention of a 2,000 km telegraph line in the Amazon connects this to the specific "Telegraph Expedition" led by Cândido Rondon. This infrastructure was crucial for both communication and the subsequent demarcation of indigenous territories.} \\
\\ 
\texttt{GUESS: The Rondon Telegraph Expedition (1907–1912).} \\
            }
        }
    }
     \\ \\

    \theutterance \stepcounter{utterance}  
    & & \multicolumn{4}{p{0.6\linewidth}}{
        \cellcolor[rgb]{0.9,0.9,0.9}{
            \makecell[{{p{\linewidth}}}]{
                \texttt{\tiny{[P1$\langle$GM]}}
                \texttt{Round 2 complete. The Detective's current guess: "The Rondon Telegraph Expedition (1907–1912).". Please write a new paragraph with fresh information.} \\
            }
        }
    }
    & \\ \\

    \theutterance \stepcounter{utterance}  
    & \multicolumn{4}{p{0.6\linewidth}}{
        \cellcolor[rgb]{0.9,0.9,0.9}{
            \makecell[{{p{\linewidth}}}]{
                \texttt{\tiny{[P1$\rangle$GM]}}
                \texttt{PARAGRAPH: The success of this arduous journey was so profound that it inspired the creation of a permanent, state{-}sponsored organization dedicated to continuing the work of mapping the frontier, training indigenous guides, and establishing a network of schools and health posts across the vast interior of the nation.} \\
            }
        }
    }
    & & \\ \\

    \theutterance \stepcounter{utterance}  
    & & \multicolumn{4}{p{0.6\linewidth}}{
        \cellcolor[rgb]{0.9,0.9,0.9}{
            \makecell[{{p{\linewidth}}}]{
                \texttt{\tiny{[GM$\rangle$P2]}}
                \texttt{{[}Sentence 3{]}} \\
\texttt{The success of this arduous journey was so profound that it inspired the creation of a permanent, state{-}sponsored organization dedicated to continuing the work of mapping the frontier, training indigenous guides, and establishing a network of schools and health posts across the vast interior of the nation.} \\
            }
        }
    }
    & \\ \\

    \theutterance \stepcounter{utterance}  
    & & & \multicolumn{4}{p{0.6\linewidth}}{
        \cellcolor[rgb]{0.9,0.9,0.9}{
            \makecell[{{p{\linewidth}}}]{
                \texttt{\tiny{[GM$\langle$P2]}}
                \texttt{ANALYSIS: The description confirms the establishment of a permanent state organization following the expedition's success. This organization focused on mapping, training indigenous guides, and building schools/health posts, which aligns perfectly with the legacy of Cândido Rondon's work.} \\
\\ 
\texttt{GUESS: The creation of the Serviço de Proteção aos Índios (SPI) in 1910.} \\
            }
        }
    }
     \\ \\

    \theutterance \stepcounter{utterance}  
    & & \multicolumn{4}{p{0.6\linewidth}}{
        \cellcolor[rgb]{0.9,0.9,0.9}{
            \makecell[{{p{\linewidth}}}]{
                \texttt{\tiny{[P1$\langle$GM]}}
                \texttt{Round 3 complete. The Detective's current guess: "The creation of the Serviço de Proteção aos Índios (SPI) in 1910.". Please write a new paragraph with fresh information.} \\
            }
        }
    }
    & \\ \\

    \theutterance \stepcounter{utterance}  
    & \multicolumn{4}{p{0.6\linewidth}}{
        \cellcolor[rgb]{0.9,0.9,0.9}{
            \makecell[{{p{\linewidth}}}]{
                \texttt{\tiny{[P1$\rangle$GM]}}
                \texttt{PARAGRAPH: This enduring institution eventually evolved into a powerful federal agency that would oversee the protection of indigenous populations for decades, fundamentally reshaping the relationship between the Brazilian state and the native peoples of the interior through a combination of assimilationist policies and administrative oversight.} \\
            }
        }
    }
    & & \\ \\

    \theutterance \stepcounter{utterance}  
    & & \multicolumn{4}{p{0.6\linewidth}}{
        \cellcolor[rgb]{0.9,0.9,0.9}{
            \makecell[{{p{\linewidth}}}]{
                \texttt{\tiny{[GM$\rangle$P2]}}
                \texttt{{[}Sentence 4{]}} \\
\texttt{This enduring institution eventually evolved into a powerful federal agency that would oversee the protection of indigenous populations for decades, fundamentally reshaping the relationship between the Brazilian state and the native peoples of the interior through a combination of assimilationist policies and administrative oversight.} \\
            }
        }
    }
    & \\ \\

    \theutterance \stepcounter{utterance}  
    & & & \multicolumn{4}{p{0.6\linewidth}}{
        \cellcolor[rgb]{0.9,0.9,0.9}{
            \makecell[{{p{\linewidth}}}]{
                \texttt{\tiny{[GM$\langle$P2]}}
                \texttt{ANALYSIS: The description confirms the evolution of the SPI into a powerful federal agency with a dual mandate of protection and assimilation. This aligns with the historical trajectory of the SPI under Rondon's influence, eventually becoming the National Indian Foundation (FUNAI).} \\
\\ 
\texttt{GUESS: The evolution of the Serviço de Proteção aos Índios (SPI) into the National Indian Foundation (FUNAI).} \\
            }
        }
    }
     \\ \\

    \theutterance \stepcounter{utterance}  
    & & \multicolumn{4}{p{0.6\linewidth}}{
        \cellcolor[rgb]{0.9,0.9,0.9}{
            \makecell[{{p{\linewidth}}}]{
                \texttt{\tiny{[P1$\langle$GM]}}
                \texttt{Round 4 complete. The Detective's current guess: "The evolution of the Serviço de Proteção aos Índios (SPI) into the National Indian Foundation (FUNAI).". Please write a new paragraph with fresh information.} \\
            }
        }
    }
    & \\ \\

    \theutterance \stepcounter{utterance}  
    & \multicolumn{4}{p{0.6\linewidth}}{
        \cellcolor[rgb]{0.9,0.9,0.9}{
            \makecell[{{p{\linewidth}}}]{
                \texttt{\tiny{[P1$\rangle$GM]}}
                \texttt{PARAGRAPH: The legacy of this initial expedition was cemented when the very explorer who led it was later honored with a prestigious national award, a recognition that highlighted how his decades of service in the Amazon had permanently altered the trajectory of Brazilian environmental and social policy.} \\
            }
        }
    }
    & & \\ \\

    \theutterance \stepcounter{utterance}  
    & & \multicolumn{4}{p{0.6\linewidth}}{
        \cellcolor[rgb]{0.9,0.9,0.9}{
            \makecell[{{p{\linewidth}}}]{
                \texttt{\tiny{[GM$\rangle$P2]}}
                \texttt{{[}Sentence 5{]}} \\
\texttt{The legacy of this initial expedition was cemented when the very explorer who led it was later honored with a prestigious national award, a recognition that highlighted how his decades of service in the Amazon had permanently altered the trajectory of Brazilian environmental and social policy.} \\
            }
        }
    }
    & \\ \\

    \theutterance \stepcounter{utterance}  
    & & & \multicolumn{4}{p{0.6\linewidth}}{
        \cellcolor[rgb]{0.9,0.9,0.9}{
            \makecell[{{p{\linewidth}}}]{
                \texttt{\tiny{[GM$\langle$P2]}}
                \texttt{ANALYSIS: The explorer receiving a prestigious national award confirms the identity of the leader as a celebrated figure. This likely refers to Rondon receiving the Legion of Honor or a Brazilian national order, cementing his status as a hero who bridged indigenous and state interests.} \\
\\ 
\texttt{GUESS: The awarding of the Legion of Honor to Cândido Rondon in 1912.} \\
            }
        }
    }
     \\ \\

    \theutterance \stepcounter{utterance}  
    & & \multicolumn{4}{p{0.6\linewidth}}{
        \cellcolor[rgb]{0.9,0.9,0.9}{
            \makecell[{{p{\linewidth}}}]{
                \texttt{\tiny{[P1$\langle$GM]}}
                \texttt{Round 5 complete. The Detective's current guess: "The awarding of the Legion of Honor to Cândido Rondon in 1912.". Please write a new paragraph with fresh information.} \\
            }
        }
    }
    & \\ \\

    \theutterance \stepcounter{utterance}  
    & \multicolumn{4}{p{0.6\linewidth}}{
        \cellcolor[rgb]{0.9,0.9,0.9}{
            \makecell[{{p{\linewidth}}}]{
                \texttt{\tiny{[P1$\rangle$GM]}}
                \texttt{PARAGRAPH: The ultimate culmination of this entire historical arc occurred in 1953, when the specific organization founded by the explorer was officially dissolved and its responsibilities were transferred to a new, independent federal entity, marking the end of an era that had begun with a single, ambitious survey mission.} \\
            }
        }
    }
    & & \\ \\

    \theutterance \stepcounter{utterance}  
    & & \multicolumn{4}{p{0.6\linewidth}}{
        \cellcolor[rgb]{0.9,0.9,0.9}{
            \makecell[{{p{\linewidth}}}]{
                \texttt{\tiny{[GM$\rangle$P2]}}
                \texttt{{[}Sentence 6{]}} \\
\texttt{The ultimate culmination of this entire historical arc occurred in 1953, when the specific organization founded by the explorer was officially dissolved and its responsibilities were transferred to a new, independent federal entity, marking the end of an era that had begun with a single, ambitious survey mission.} \\
            }
        }
    }
    & \\ \\

    \theutterance \stepcounter{utterance}  
    & & & \multicolumn{4}{p{0.6\linewidth}}{
        \cellcolor[rgb]{0.9,0.9,0.9}{
            \makecell[{{p{\linewidth}}}]{
                \texttt{\tiny{[GM$\langle$P2]}}
                \texttt{ANALYSIS: The dissolution of the organization in 1953 marks the end of the SPI's direct administration. This event transferred indigenous affairs to a new entity, likely FUNAI, concluding the specific administrative era started by Rondon's expedition.} \\
\\ 
\texttt{GUESS: The dissolution of the Serviço de Proteção aos Índios (SPI) in 1953.} \\
            }
        }
    }
     \\ \\

    \theutterance \stepcounter{utterance}  
    & & \multicolumn{4}{p{0.6\linewidth}}{
        \cellcolor[rgb]{0.9,0.9,0.9}{
            \makecell[{{p{\linewidth}}}]{
                \texttt{\tiny{[P1$\langle$GM]}}
                \texttt{Round 6 complete. The Detective's current guess: "The dissolution of the Serviço de Proteção aos Índios (SPI) in 1953.". Please write a new paragraph with fresh information.} \\
            }
        }
    }
    & \\ \\

    \theutterance \stepcounter{utterance}  
    & \multicolumn{4}{p{0.6\linewidth}}{
        \cellcolor[rgb]{0.9,0.9,0.9}{
            \makecell[{{p{\linewidth}}}]{
                \texttt{\tiny{[P1$\rangle$GM]}}
                \texttt{PARAGRAPH: The final chapter of this long saga was written when the original explorer, who had passed away years prior, was posthumously awarded the highest military decoration in the nation, a tribute that acknowledged his lifelong dedication to the very borders and peoples that the entire complex of missions, telegraph lines, and agencies had been built to protect and define.} \\
            }
        }
    }
    & & \\ \\

    \theutterance \stepcounter{utterance}  
    & & \multicolumn{4}{p{0.6\linewidth}}{
        \cellcolor[rgb]{0.9,0.9,0.9}{
            \makecell[{{p{\linewidth}}}]{
                \texttt{\tiny{[GM$\rangle$P2]}}
                \texttt{{[}Sentence 7{]}} \\
\texttt{The final chapter of this long saga was written when the original explorer, who had passed away years prior, was posthumously awarded the highest military decoration in the nation, a tribute that acknowledged his lifelong dedication to the very borders and peoples that the entire complex of missions, telegraph lines, and agencies had been built to protect and define.} \\
            }
        }
    }
    & \\ \\

    \theutterance \stepcounter{utterance}  
    & & & \multicolumn{4}{p{0.6\linewidth}}{
        \cellcolor[rgb]{0.9,0.9,0.9}{
            \makecell[{{p{\linewidth}}}]{
                \texttt{\tiny{[GM$\langle$P2]}}
                \texttt{ANALYSIS: The explorer, Cândido Rondon, died in 1958. The highest military decoration in Brazil is the Order of the Southern Cross. This sentence describes the posthumous awarding of this specific order to him, capping his legacy of border demarcation and indigenous protection.} \\
\\ 
\texttt{GUESS: The posthumous awarding of the Order of the Southern Cross to Cândido Rondon in 1958.} \\
            }
        }
    }
     \\ \\

    \theutterance \stepcounter{utterance}  
    & & \multicolumn{4}{p{0.6\linewidth}}{
        \cellcolor[rgb]{0.9,0.9,0.9}{
            \makecell[{{p{\linewidth}}}]{
                \texttt{\tiny{[P1$\langle$GM]}}
                \texttt{Round 7 complete. The Detective's current guess: "The posthumous awarding of the Order of the Southern Cross to Cândido Rondon in 1958.". Please write a new paragraph with fresh information.} \\
            }
        }
    }
    & \\ \\

    \theutterance \stepcounter{utterance}  
    & \multicolumn{4}{p{0.6\linewidth}}{
        \cellcolor[rgb]{0.9,0.9,0.9}{
            \makecell[{{p{\linewidth}}}]{
                \texttt{\tiny{[P1$\rangle$GM]}}
                \texttt{PARAGRAPH: The entire sequence of events you have been piecing together began with a specific presidential decree issued in 1907, which formally authorized the expedition that would go on to build the telegraph line, establish the agency, and ultimately define the modern concept of indigenous protection in Brazil.} \\
            }
        }
    }
    & & \\ \\

    \theutterance \stepcounter{utterance}  
    & & \multicolumn{4}{p{0.6\linewidth}}{
        \cellcolor[rgb]{0.9,0.9,0.9}{
            \makecell[{{p{\linewidth}}}]{
                \texttt{\tiny{[GM$\rangle$P2]}}
                \texttt{{[}Sentence 8{]}} \\
\texttt{The entire sequence of events you have been piecing together began with a specific presidential decree issued in 1907, which formally authorized the expedition that would go on to build the telegraph line, establish the agency, and ultimately define the modern concept of indigenous protection in Brazil.} \\
            }
        }
    }
    & \\ \\

    \theutterance \stepcounter{utterance}  
    & & & \multicolumn{4}{p{0.6\linewidth}}{
        \cellcolor[rgb]{0.9,0.9,0.9}{
            \makecell[{{p{\linewidth}}}]{
                \texttt{\tiny{[GM$\langle$P2]}}
                \texttt{ANALYSIS: This sentence identifies the specific legal origin of the entire saga: the 1907 presidential decree that authorized the expedition. This confirms the starting point of the historical arc involving Rondon, the telegraph, and the SPI.} \\
\\ 
\texttt{GUESS: The 1907 Presidential Decree authorizing the Rondon Telegraph Expedition.} \\
            }
        }
    }
     \\ \\

    \theutterance \stepcounter{utterance}  
    & & \multicolumn{4}{p{0.6\linewidth}}{
        \cellcolor[rgb]{0.9,0.9,0.9}{
            \makecell[{{p{\linewidth}}}]{
                \texttt{\tiny{[P1$\langle$GM]}}
                \texttt{Round 8 complete. The Detective's current guess: "The 1907 Presidential Decree authorizing the Rondon Telegraph Expedition.". Please write a new paragraph with fresh information.} \\
            }
        }
    }
    & \\ \\

    \theutterance \stepcounter{utterance}  
    & \multicolumn{4}{p{0.6\linewidth}}{
        \cellcolor[rgb]{0.9,0.9,0.9}{
            \makecell[{{p{\linewidth}}}]{
                \texttt{\tiny{[P1$\rangle$GM]}}
                \texttt{PARAGRAPH: This pivotal moment in history was not merely a surveying mission but the genesis of a unique partnership between the Brazilian government and indigenous communities, creating a model where native guides were employed as essential partners in the nation's expansion and infrastructure development rather than as subjects to be displaced.} \\
            }
        }
    }
    & & \\ \\

    \theutterance \stepcounter{utterance}  
    & & \multicolumn{4}{p{0.6\linewidth}}{
        \cellcolor[rgb]{0.9,0.9,0.9}{
            \makecell[{{p{\linewidth}}}]{
                \texttt{\tiny{[GM$\rangle$P2]}}
                \texttt{{[}Sentence 9{]}} \\
\texttt{This pivotal moment in history was not merely a surveying mission but the genesis of a unique partnership between the Brazilian government and indigenous communities, creating a model where native guides were employed as essential partners in the nation's expansion and infrastructure development rather than as subjects to be displaced.} \\
            }
        }
    }
    & \\ \\

    \theutterance \stepcounter{utterance}  
    & & & \multicolumn{4}{p{0.6\linewidth}}{
        \cellcolor[rgb]{0.9,0.9,0.9}{
            \makecell[{{p{\linewidth}}}]{
                \texttt{\tiny{[GM$\langle$P2]}}
                \texttt{ANALYSIS: This sentence emphasizes the unique "partnership" model where indigenous people were employed as guides rather than displaced. This specific dynamic defines the core philosophy of the Rondon Expedition and the subsequent SPI, distinguishing it from other colonial expansion efforts.} \\
\\ 
\texttt{GUESS: The Rondon Expedition (1907–1912) and the subsequent establishment of the Serviço de Proteção aos Índios (SPI).} \\
            }
        }
    }
     \\ \\

    \theutterance \stepcounter{utterance}  
    & & \multicolumn{4}{p{0.6\linewidth}}{
        \cellcolor[rgb]{0.9,0.9,0.9}{
            \makecell[{{p{\linewidth}}}]{
                \texttt{\tiny{[P1$\langle$GM]}}
                \texttt{Round 9 complete. The Detective's current guess: "The Rondon Expedition (1907–1912) and the subsequent establishment of the Serviço de Proteção aos Índios (SPI).". Please write a new paragraph with fresh information.} \\
            }
        }
    }
    & \\ \\

    \theutterance \stepcounter{utterance}  
    & \multicolumn{4}{p{0.6\linewidth}}{
        \cellcolor[rgb]{0.9,0.9,0.9}{
            \makecell[{{p{\linewidth}}}]{
                \texttt{\tiny{[P1$\rangle$GM]}}
                \texttt{PARAGRAPH: The comprehensive legacy of this entire historical arc is best understood as the foundational framework for Brazil's first major conservation initiative, which successfully merged the goals of national border demarcation with the protection of indigenous territories, setting a precedent that would influence environmental policy for generations to come.} \\
            }
        }
    }
    & & \\ \\

    \theutterance \stepcounter{utterance}  
    & & \multicolumn{4}{p{0.6\linewidth}}{
        \cellcolor[rgb]{0.9,0.9,0.9}{
            \makecell[{{p{\linewidth}}}]{
                \texttt{\tiny{[GM$\rangle$P2]}}
                \texttt{{[}Sentence 10{]}} \\
\texttt{The comprehensive legacy of this entire historical arc is best understood as the foundational framework for Brazil's first major conservation initiative, which successfully merged the goals of national border demarcation with the protection of indigenous territories, setting a precedent that would influence environmental policy for generations to come.} \\
            }
        }
    }
    & \\ \\

    \theutterance \stepcounter{utterance}  
    & & & \multicolumn{4}{p{0.6\linewidth}}{
        \cellcolor[rgb]{0.9,0.9,0.9}{
            \makecell[{{p{\linewidth}}}]{
                \texttt{\tiny{[GM$\langle$P2]}}
                \texttt{ANALYSIS: This final sentence synthesizes the entire narrative: the 1907 decree, the telegraph expedition, the SPI, and the unique partnership model. It identifies the overarching historical significance as the creation of Brazil's first major conservation initiative that merged border demarcation with indigenous protection.} \\
\\ 
\texttt{GUESS: The creation of the Brazilian Amazon Conservation Model (specifically the legacy of the Rondon Expedition and the SPI).} \\
            }
        }
    }
     \\ \\

    \theutterance \stepcounter{utterance}  
    & & & \multicolumn{2}{p{0.3\linewidth}}{
        \cellcolor[rgb]{0.9,0.9,0.9}{
            \makecell[{{p{\linewidth}}}]{
                \texttt{\tiny{[GM$|$GM]}}
                \texttt{10} \\
            }
        }
    }
    & & \\ \\

\end{supertabular}
}

\end{document}
